
\chapter{Guidance design}
\label{ch:guide}

The previous chapter specified the scenario to be generated through an initial
atmospheric context, a target mask \(m\), and a target trajectory
\(y^\star_{1:N}\). We now address the second component of the framework:
how the generative forecast can be steered towards this prescribed target.

We formulate this as a \emph{target-driven guidance} problem. At each weather
step \(n\), the target \(y_n^\star\) is fixed before generation begins. During
the internal residual flow, the current clean-state estimate is repeatedly
compared with this target, and the resulting discrepancy is used to modify the
generative dynamics. The central problem is therefore not to choose a guidance
strength and observe the resulting scenario, but rather to determine the
guidance required to realize a scenario that has been specified in advance.

We first introduce the common guidance formulation used throughout this
chapter. We then develop three strategies for realizing the prescribed target:
\textsc{Close-Gap}, which locally controls the remaining target gap;
\textsc{Optimize-Schedule}, which jointly optimizes the interventions across
the flow; and \textsc{Optimize-Gain}, which fixes the relative flow-time
profile and optimizes a single guidance gain. Finally, we introduce Universal
Guidance and adversarial weather forecasting as baseline approaches.

\section{Target-driven guidance}
\label{sec:target-driven-guidance}

Guidance operates independently at every weather step of the autoregressive
forecast. For notational simplicity, we denote by
\(\hat{x}_{n,t}\) the clean-state estimate of the weather state being generated
at weather step \(n\), evaluated at internal flow time \(t\). Thus, \(n\)
continues to index weather time, while \(t\) refers exclusively to the
generative flow.

Although the guidance problem is solved separately at each weather step, its
effects are not independent across weather time. Once the guided state
\(x_n^{\mathrm{gui}}\) has been generated, it becomes part of the conditioning
context used to generate the following state. Guidance applied at one weather
step can therefore propagate through the subsequent autoregressive forecast.

\subsection{Guidance objective}
\label{sec:guidance-objective}

At weather step \(n\), the target construction introduced in
\Cref{ch:specify} provides a scalar target value \(y_n^\star\). Given the mask
\(m\), we measure the discrepancy between the current clean-state estimate and
this target by

\begin{equation}
    \xi_{n,t}
    :=
    \mathcal{M}_m(\hat{x}_{n,t}) - y_n^\star ,
    \label{eq:guidance-gap}
\end{equation}

where \(\mathcal{M}_m\) denotes the normalized masked average defined in
\Cref{ch:specify}. We refer to \(\xi_{n,t}\) as the \emph{target gap}.

The corresponding guidance loss is the squared target gap,

\begin{equation}
    \mathcal{L}_{n,t}
    :=
    \xi_{n,t}^{\,2}
    =
    \left(
        \mathcal{M}_m(\hat{x}_{n,t}) - y_n^\star
    \right)^2 .
    \label{eq:guidance-loss}
\end{equation}

This formulation separates the specification of the desired scenario from the
mechanism used to realize it. In the controlled experiments of
\Cref{ch:specify}, for example,

\begin{equation}
    y_n^\star
    =
    (1+\rho_n)\,
    \mathcal{M}_m(x_n^{\mathrm{ung}}),
    \label{eq:controlled-guidance-target}
\end{equation}

whereas for historical storyline generation \(y_n^\star\) is obtained from a
selected historical trajectory. The guidance algorithms developed below are
unchanged by this choice: once \(y_n^\star\) has been specified, they operate
only on the resulting target gap.

After completing the internal flow, the generated weather state is denoted by
\(x_n^{\mathrm{gui}}\). Its final realized target gap is therefore

\begin{equation}
    \xi_n
    :=
    \mathcal{M}_m(x_n^{\mathrm{gui}}) - y_n^\star
    =
    \xi_{n,T}.
    \label{eq:final-guidance-gap}
\end{equation}

The objective of the guidance methods in this chapter is to reduce
\(\lvert\xi_n\rvert\), ideally reaching \(\xi_n=0\). Importantly, this does
not imply that arbitrary changes to the atmospheric state are considered
acceptable. Reaching the prescribed target describes only whether guidance
has succeeded with respect to the specified event; the meteorological
plausibility of the resulting forecast is evaluated separately.

\subsection{Gradient-based flow guidance}
\label{sec:gradient-flow-guidance}

As introduced in the background chapter, loss-based guidance modifies the
generative velocity using the gradient of a differentiable objective with
respect to the current noisy state. At weather step \(n\) and flow time \(t\),
we compute

\begin{equation}
    g_{n,t}
    :=
    \nabla_{z_{n,t}}
    \mathcal{L}_{n,t}.
    \label{eq:guidance-gradient}
\end{equation}

For the squared target-gap loss in \Cref{eq:guidance-loss}, this becomes

\begin{equation}
    g_{n,t}
    =
    2\xi_{n,t}\,
    \nabla_{z_{n,t}}
    \mathcal{M}_m(\hat{x}_{n,t}).
    \label{eq:guidance-gradient-expanded}
\end{equation}

Although the objective itself is defined only through the masked quantity
\(\mathcal{M}_m\), the gradient is taken with respect to the complete noisy
residual \(z_{n,t}\). It therefore propagates through the mapping from the
noisy state to the clean-state estimate and can contain non-zero components
outside the explicitly targeted variable and region. The intervention is thus
not obtained by directly adding the mask to the generated weather field.
Instead, the gradient describes the local change of the noisy state that
modifies the masked quantity through the learned generative model.

Using the notation introduced in the background chapter, the guided velocity
can be written as

\begin{equation}
    \tilde{u}_{n,t}
    =
    u_{n,t}
    -
    \lambda_{n,t} g_{n,t},
    \label{eq:target-driven-guided-velocity}
\end{equation}

followed by the usual Euler update

\begin{equation}
    z_{n,t+1}
    =
    z_{n,t}
    +
    h_t \tilde{u}_{n,t}.
    \label{eq:guided-euler-update}
\end{equation}

It is useful to separate the direction of the guidance gradient from its
magnitude. Defining

\begin{equation}
    \bar{g}_{n,t}
    :=
    \frac{g_{n,t}}{\lVert g_{n,t}\rVert_2},
    \label{eq:unit-guidance-gradient}
\end{equation}

the same velocity modification can equivalently be expressed as

\begin{equation}
    \tilde{u}_{n,t}
    =
    u_{n,t}
    -
    q_{n,t}\bar{g}_{n,t},
    \qquad
    q_{n,t}
    :=
    \lambda_{n,t}\lVert g_{n,t}\rVert_2 .
    \label{eq:unit-gradient-guidance}
\end{equation}

This representation will be convenient when comparing the methods developed
below. Some methods directly control the multiplier of the raw gradient,
whereas others naturally control the magnitude of the resulting intervention.
Equations~\eqref{eq:target-driven-guided-velocity} and
\eqref{eq:unit-gradient-guidance} are equivalent representations of the same
update.

\subsection{Guidance-strength parameterization}
\label{sec:guidance-strength-parameterization}

We decompose the raw-gradient multiplier introduced in
\Cref{eq:target-driven-guided-velocity} as

\begin{equation}
    \lambda_{n,t}
    =
    w_n\,\gamma_t\,\nu_{n,t}.
    \label{eq:guidance-strength-decomposition}
\end{equation}

The three terms play distinct roles. The scalar \(w_n\) controls the overall
guidance magnitude at weather step \(n\). The profile \(\gamma_t\) controls
how this magnitude is distributed over the internal flow. Finally,
\(\nu_{n,t}\) represents an optional method-specific normalization or scaling
factor.

We take \(\gamma_t\) to be a non-negative normalized profile with maximum
value one. Only its relative variation across flow time is relevant: any
positive multiplicative rescaling of \(\gamma_t\) can be absorbed into
\(w_n\). The concrete profile families considered in this thesis are therefore
an experimental design choice and are introduced later together with the
experimental setup.

The decomposition in \Cref{eq:guidance-strength-decomposition} also clarifies
the distinction between the methods developed below. In
\textsc{Optimize-Gain}, the flow-time profile \(\gamma_t\) is fixed and only
the scalar \(w_n\) is optimized. In \textsc{Optimize-Schedule}, the
interventions at individual flow steps are instead optimized jointly. In
\textsc{Close-Gap}, the profile has a different interpretation altogether:
rather than prescribing the relative strength of guidance, it prescribes the
desired evolution of the remaining target gap. We make this distinction
explicit when deriving each method.

Finally, we do not apply guidance at the last Euler step of the residual flow.
For ArchesWeather, the final step size satisfies \(h_t=s_t\). Since

\begin{equation}
    u_{n,t}
    =
    \frac{\hat r_{n,t}-z_{n,t}}{s_t},
\end{equation}

the unguided final Euler update gives

\begin{equation}
    z_{n,t+1}
    =
    z_{n,t}
    +
    s_t
    \frac{\hat r_{n,t}-z_{n,t}}{s_t}
    =
    \hat r_{n,t}.
    \label{eq:final-flow-snap}
\end{equation}

A guidance correction applied at this point would therefore directly alter
the final residual without any subsequent model evaluation in which the
perturbed state could be processed. We consequently set the guidance
magnitude to zero at the final flow step for all methods considered in this
chapter.

\section{Target-realization methods}
\label{sec:target-realization-methods}

The formulation above leaves open how the guidance magnitude should be chosen
in order to realize the prescribed target. 
---this is correct, but it's more that we develop them, although in a way they follow the theoeritical framework we have developed
We investigate three approaches to
this problem. They differ primarily in the number of guidance degrees of
freedom that are controlled.

\textsc{Close-Gap} determines a separate correction at every flow step from
the current target gap. \textsc{Optimize-Schedule} instead treats the complete
sequence of flow-time interventions as an optimization variable.
\textsc{Optimize-Gain} reduces this problem to a single scalar by fixing the
relative flow-time profile and optimizing only its overall magnitude.

All three methods act on the same loss introduced in
\Cref{eq:guidance-loss} and are applied independently at every weather step
\(n\).


\subsection{\textsc{Close-Gap}}
\label{sec:close-gap}

Since the desired target is known in advance, a direct strategy is to prescribe
not only the final target gap ($\approx0$), but also how this gap should evolve throughout
the flow. Let

\begin{equation}
    \bar{\xi}_{n,t}
    :=
    \gamma_t^{\mathrm{gap}}\,\xi_{n,0},
    \label{eq:close-gap-path}
\end{equation}

denote the prescribed remaining gap at flow step $t$, where
$\gamma_t^{\mathrm{gap}}$ controls the fraction of the initial gap that should
remain. A decreasing profile therefore progressively closes the gap, with
$\gamma_t^{\mathrm{gap}}\rightarrow 0$ corresponding to complete closure.

At each flow step, \textsc{Close-Gap} computes the local displacement of the
noisy state required to move the current gap $\xi_{n,t}$ towards
$\bar{\xi}_{n,t}$. For a small displacement $\Delta z_{n,t}$, we approximate the target gap
using a first-order Taylor expansion around the current noisy state:

\begin{equation}
    \xi_{n,t}(z_{n,t}+\Delta z_{n,t})
    \approx
    \xi_{n,t}(z_{n,t})
    +
    \left\langle
        \nabla_{z_{n,t}}\xi_{n,t},
        \Delta z_{n,t}
    \right\rangle.
\end{equation}

Requiring the corrected gap to equal the prescribed value gives the local
constraint

\begin{equation}
    \left\langle
        \nabla_{z_{n,t}}\xi_{n,t},
        \Delta z_{n,t}
    \right\rangle
    =
    \bar{\xi}_{n,t}-\xi_{n,t}.
    \label{eq:close-gap-constraint}
\end{equation}

Among all displacements satisfying this first-order constraint, the
minimum-norm solution is aligned with the gradient of the gap,

\begin{equation}
    \Delta z_{n,t}^{\star}
    =
    \frac{
        \bar{\xi}_{n,t}-\xi_{n,t}
    }{
        \left\|
            \nabla_{z_{n,t}}\xi_{n,t}
        \right\|_2^2
    }
    \nabla_{z_{n,t}}\xi_{n,t}.
    \label{eq:close-gap-minimum-shift}
\end{equation}

The guidance gradient introduced above is

\begin{equation}
    g_{n,t}
    =
    \nabla_{z_{n,t}}\xi_{n,t}^{\,2}
    =
    2\xi_{n,t}
    \nabla_{z_{n,t}}\xi_{n,t}.
    \label{eq:close-gap-gradient}
\end{equation}

Substituting this relation into
\Cref{eq:close-gap-minimum-shift} gives the desired guidance-induced
displacement directly in terms of $g_{n,t}$,

\begin{equation}
    \Delta z_{n,t}^{\star}
    =
    -
    \frac{
        2\xi_{n,t}
        \left(
            \xi_{n,t}-\bar{\xi}_{n,t}
        \right)
    }{
        \lVert g_{n,t}\rVert_2^2
    }
    g_{n,t}.
    \label{eq:close-gap-shift}
\end{equation}

On the other hand, modifying the velocity according to

\begin{equation}
    \tilde{u}_{n,t}
    =
    u_{n,t}
    -
    \lambda_{n,t}g_{n,t}
\end{equation}

introduces the following guidance contribution to the Euler step:

\begin{equation}
    \Delta z_{n,t}^{\mathrm{gui}}
    =
    -h_t\lambda_{n,t}g_{n,t}.
    \label{eq:close-gap-guidance-shift}
\end{equation}

Requiring the guidance contribution to the Euler step to realize the desired
local displacement in \Cref{eq:close-gap-shift} yields the closed-form
guidance multiplier

\begin{equation}
    \lambda_{n,t}^{\mathrm{CG}}
    =
    \frac{
        2\xi_{n,t}
        \left(
            \xi_{n,t}-\bar{\xi}_{n,t}
        \right)
    }{
        h_t
        \lVert g_{n,t}\rVert_2^2
    }.
    \label{eq:close-gap-lambda}
\end{equation}

When the current gap already equals the prescribed gap, the required
correction is zero. Otherwise, the multiplier is recomputed at the following
flow step using the newly obtained state.

The method is therefore greedy: at every flow step it computes the
minimum-norm local correction predicted to move the current estimate onto the
prescribed gap trajectory. This correction is derived from a first-order
approximation and accounts only for the guidance-induced displacement, while
the learned velocity field simultaneously advances the noisy state. The actual
gap may therefore deviate from the prescribed path, in which case the
difference is observed and corrected again at the next flow step.

Unlike the optimization-based methods introduced below, \textsc{Close-Gap}
requires no repeated generation of the same weather state. The guidance
strength is determined online during a single flow. Its main limitation is
that each intervention is chosen locally rather than according to its effect
on the final generated state.

\subsection{\textsc{Optimize-Schedule}}
\label{sec:optimize-schedule}

\textsc{Close-Gap} prescribes the evolution of the target gap and derives a
local intervention at each flow step. A more general alternative is to optimize
the interventions across the complete flow jointly, based on their effect on
the final generated state.

We parameterize the intervention at each flow step through its contribution to
the Euler update,

\begin{equation}
    k_{n,t}
    :=
    h_t\lambda_{n,t},
    \label{eq:guidance-kick}
\end{equation}

such that the guided Euler step becomes

\begin{equation}
    z_{n,t+1}
    =
    z_{n,t}
    +
    h_t u_{n,t}
    -
    k_{n,t}g_{n,t}.
    \label{eq:kick-euler-update}
\end{equation}

We refer to $k_{n,t}g_{n,t}$ as the \emph{guidance kick}. This
parameterization is convenient because the kick corresponds directly to the
guidance-induced displacement of the noisy state during one Euler step. The
magnitude of the intervention can therefore be regularized directly in the
same space in which it is applied.

Since the final Euler step is not guided, we optimize the vector

\begin{equation}
    \mathbf{k}_n
    =
    \left(
        k_{n,0},\ldots,k_{n,T-2}
    \right)
\end{equation}

according to

\begin{equation}
    J_n(\mathbf{k}_n)
    =
    \xi_n^2
    +
    \kappa
    \sum_{t=0}^{T-2}
    \left\|
        k_{n,t}g_{n,t}
    \right\|_2^2.
    \label{eq:optimize-schedule-objective}
\end{equation}

The first term measures the final target gap, while the second penalizes the
total squared magnitude of the guidance-induced state displacements.
The coefficient $\kappa\geq0$ therefore controls the trade-off between target
realization and intervention magnitude.

This objective mirrors the terminal-loss and control-penalty structure
introduced in \Cref{sec:guidance-as-control}. Here, however, the control
direction is restricted to the guidance gradient $g_{n,t}$, while its magnitude
is optimized independently at each flow step. The regularization is defined
directly on the discrete state displacements applied by the numerical sampler.

\paragraph{Gradient with respect to the schedule.}

Optimizing $\mathbf{k}_n$ requires differentiating the final target loss
through the sequence of Euler updates. Following the decomposed
backpropagation strategy of \textcite{liu2023flowgrad}, we propagate the
terminal gradient backwards through the discretized flow using
vector--Jacobian products, without retaining the complete unrolled
computation graph.

We express this backward recursion as a discrete adjoint. Let

\begin{equation}
    p_{n,T}
    :=
    \nabla_{z_{n,T}}\xi_n^2
\end{equation}

denote the terminal adjoint. In addition, we introduce a
frozen-guidance-gradient approximation: the guidance gradient $g_{n,t}$ is
recomputed during each forward evaluation of the schedule, but treated as
fixed during the corresponding backward pass. Under this approximation, the
Jacobian of the Euler update with respect to the noisy state is

\begin{equation}
    \frac{\partial z_{n,t+1}}
         {\partial z_{n,t}}
    \approx
    I
    +
    h_t
    \frac{\partial u_{n,t}}
         {\partial z_{n,t}},
\end{equation}

and the adjoint is propagated backwards according to

\begin{equation}
    p_{n,t}
    \approx
    p_{n,t+1}
    +
    h_t
    \left(
        \frac{\partial u_{n,t}}
             {\partial z_{n,t}}
    \right)^{\!\top}
    p_{n,t+1}.
    \label{eq:frozen-gradient-adjoint}
\end{equation}
The adjoint $p_{n,t+1}$ describes how a perturbation of the state
$z_{n,t+1}$ affects the final target loss. Since the kick $k_{n,t}$
enters the Euler update in \Cref{eq:kick-euler-update} as
$-k_{n,t}g_{n,t}$, the derivative of the final target loss with respect
to the kick is

\begin{equation}
    \frac{\partial \xi_n^2}
         {\partial k_{n,t}}
    \approx
    -
    \left\langle
        p_{n,t+1},
        g_{n,t}
    \right\rangle.
    \label{eq:kick-target-gradient}
\end{equation}

Including the intervention penalty in
\Cref{eq:optimize-schedule-objective} gives

\begin{equation}
    \frac{\partial J_n}
         {\partial k_{n,t}}
    \approx
    -
    \left\langle
        p_{n,t+1},
        g_{n,t}
    \right\rangle
    +
    2\kappa
    k_{n,t}
    \lVert g_{n,t}\rVert_2^2.
    \label{eq:kick-objective-gradient}
\end{equation}

These gradients are used to jointly optimize the non-negative kick trajectory
$\mathbf{k}_n$. The optimizer, initialization, stopping criteria, and
regularization strength are specified with the experimental setup.

Compared with \textsc{Close-Gap}, \textsc{Optimize-Schedule} accounts for how
interventions applied at different flow steps influence the final target gap.
This comes at a higher computational cost: the method introduces $T-1$
optimization variables and requires repeated forward and backward evaluations
of the flow for each generated weather state.

\subsection{\textsc{Optimize-Gain}}
\label{sec:optimize-gain}

The flexibility of \textsc{Optimize-Schedule} comes at a substantial
optimization cost. We therefore consider an intermediate formulation in which
the relative distribution of guidance over flow time is fixed, while its
overall magnitude remains adaptive.

Let $\gamma_t$ be a prescribed normalized guidance profile. We set

\begin{equation}
    \lambda_{n,t}
    =
    w_n\gamma_t,
    \label{eq:optimize-gain-lambda}
\end{equation}

such that all flow-time guidance strengths are determined by the single scalar
$w_n\geq0$. Since the target is specified separately at every weather step,
$w_n$ is also optimized independently for each weather step.

For a fixed initial noise realization, running the complete guided flow defines
the final target loss as a scalar function

\begin{equation}
    \mathcal{L}_n(w_n)
    :=
    \xi_n(w_n)^2.
    \label{eq:gain-loss}
\end{equation}

We then seek a stationary point

\begin{equation}
    \frac{\mathrm{d}\mathcal{L}_n}
         {\mathrm{d}w_n}
    =
    0.
    \label{eq:gain-stationarity}
\end{equation}

If the prescribed target can be reached exactly, any $w_n^\star$ satisfying
$\xi_n(w_n^\star)=0$ also satisfies \Cref{eq:gain-stationarity}. If exact
closure is not reached, we retain the evaluated value of $w_n$ that produces
the smallest remaining squared gap.

Using \Cref{eq:optimize-gain-lambda}, the derivative follows directly from the
chain rule,

\begin{equation}
    \frac{\mathrm{d}\mathcal{L}_n}
         {\mathrm{d}w_n}
    =
    \sum_{t=0}^{T-2}
    \frac{\partial \mathcal{L}_n}
         {\partial \lambda_{n,t}}
    \frac{\partial \lambda_{n,t}}
         {\partial w_n}
    =
    \sum_{t=0}^{T-2}
    \gamma_t
    \frac{\partial \mathcal{L}_n}
         {\partial \lambda_{n,t}}.
    \label{eq:gain-chain-rule}
\end{equation}

The individual derivatives are evaluated using the frozen-guidance-gradient
adjoint introduced in \Cref{eq:frozen-gradient-adjoint}. Since
$\lambda_{n,t}$ enters the Euler update through
$-h_t\lambda_{n,t}g_{n,t}$,

\begin{equation}
    \frac{\partial \mathcal{L}_n}
         {\partial \lambda_{n,t}}
    \approx
    -h_t
    \left\langle
        p_{n,t+1},
        g_{n,t}
    \right\rangle.
    \label{eq:gain-lambda-gradient}
\end{equation}

Because the problem is one-dimensional, we solve
\Cref{eq:gain-stationarity} using the secant method. Given two previous
evaluations $w_n^{(j-1)}$ and $w_n^{(j)}$, the next estimate is

\begin{equation}
    w_n^{(j+1)}
    =
    w_n^{(j)}
    -
    \frac{
        \mathcal{L}_n'(w_n^{(j)})
        \left(
            w_n^{(j)}-w_n^{(j-1)}
        \right)
    }{
        \mathcal{L}_n'(w_n^{(j)})
        -
        \mathcal{L}_n'(w_n^{(j-1)})
    }.
    \label{eq:secant-guidance-gain}
\end{equation}

All evaluations use the same initial noise realization. Consequently,
$\mathcal{L}_n(w_n)$ changes only as a result of the guidance gain rather than
because of stochastic differences between repeated samples.

\textsc{Optimize-Gain} can therefore be interpreted as a restricted version
of \textsc{Optimize-Schedule}. Rather than optimizing the intervention
magnitude independently at every flow step, it restricts the guidance schedule
to the one-dimensional family

\begin{equation}
    \boldsymbol{\lambda}_n
    =
    w_n
    \boldsymbol{\gamma}.
\end{equation}

This removes most of the optimization degrees of freedom while retaining the
central target-driven property of the method: the guidance magnitude is not
chosen beforehand, but is adapted to the prescribed target.

\section{Baseline methods}
\label{sec:guidance-baselines}

The three methods above are developed specifically around the target-driven
formulation of this thesis. We compare them with two conceptually different
ways of producing a targeted intervention: Universal Guidance and adversarial
weather forecasting.

\subsection{Universal Guidance}
\label{sec:universal-guidance-baseline}

Universal Guidance was introduced as a general mechanism for controlling
diffusion models through arbitrary differentiable objectives
\parencite{bansal2023universal}. In particular, its backward-guidance
formulation first determines a change in the predicted clean sample that better
satisfies the desired condition and subsequently maps this change back to the
current noisy state.

For the objective considered in this thesis, this clean-space optimization has
a particularly simple solution. Since the masked functional
\(\mathcal{M}_m\) is linear, we seek the smallest clean-state perturbation
\(\Delta x_{n,t}\) satisfying

\begin{equation}
    \mathcal{M}_m
    \left(
        \hat{x}_{n,t}+\Delta x_{n,t}
    \right)
    =
    y_n^\star.
\end{equation}

Equivalently,

\begin{equation}
    \left\langle
        m,\Delta x_{n,t}
    \right\rangle
    =
    -\xi_{n,t}.
\end{equation}

The minimum-\(\ell_2\)-norm solution is

\begin{equation}
    \Delta x_{n,t}^{\mathrm{UG}}
    =
    -\xi_{n,t}
    \frac{m}{\lVert m\rVert_2^2}.
    \label{eq:ug-clean-shift}
\end{equation}

Thus, unlike the gradient-based methods above, no iterative clean-space
optimization is required for our particular loss.

To translate this perturbation back to flow space, let \(D\) denote the
diagonal scaling that maps the normalized residual representation used by
ArchesWeather to physical state units. The clean-state estimate can then be
written schematically as

\begin{equation}
    \hat{x}_{n,t}
    =
    \hat{x}^{\mathrm{det}}_n
    +
    D
    \left(
        z_{n,t}
        +
        s_tu_{n,t}
    \right).
    \label{eq:ug-clean-affine}
\end{equation}

The clean-space change in \Cref{eq:ug-clean-shift} therefore corresponds to the
velocity correction

\begin{equation}
    \Delta u_{n,t}^{\mathrm{UG}}
    =
    \frac{1}{s_t}
    D^{-1}
    \Delta x_{n,t}^{\mathrm{UG}}.
    \label{eq:ug-velocity-shift}
\end{equation}

We apply this correction according to the same normalized flow-time profile
\(\gamma_t\) used for \textsc{Optimize-Gain},

\begin{equation}
    \tilde{u}_{n,t}
    =
    u_{n,t}
    +
    \gamma_t
    \Delta u_{n,t}^{\mathrm{UG}}.
    \label{eq:ug-guided-velocity}
\end{equation}

This provides a useful baseline because the desired clean-state intervention is
constructed directly, rather than inferred from an optimized guidance
strength. The temporal placement of guidance can nevertheless be kept
identical to \textsc{Optimize-Gain}, allowing the effect of the intervention
mechanism itself to be compared under the same flow-time profile.

A special case occurs when \(\gamma_t\) is non-zero only at the beginning of
the flow. Universal Guidance then applies a single targeted perturbation to the
early noisy generative state, after which the remaining flow evolves without
further intervention. This can be interpreted as an adversarial-style
perturbation in the internal generative state space.

\subsection{Adversarial weather forecasting}
\label{sec:aowf-baseline}

Our second baseline acts outside the generative flow. Adversarial Observations
in Weather Forecasting (AOWF) directly optimizes perturbations to atmospheric
states in order to manipulate a subsequent data-driven forecast
\parencite{imgrund2025adversarial}. A related adversarial formulation is used
as the baseline for guided precipitation interventions by
\textcite{ueyama2026guided}.

We adapt this idea to the same target-realization task considered throughout
this chapter. When generating weather state \(x_n\), ArchesWeather is
conditioned on the preceding atmospheric states. Following the one-state
intervention setting used for comparison, we perturb the most recent
conditioning state,

\begin{equation}
    x_{n-1}
    \longrightarrow
    x_{n-1}+\delta_{n-1}.
\end{equation}

Let

\begin{equation}
    F_n
    \left(
        x_{n-2},
        x_{n-1};
        \varepsilon_n
    \right)
\end{equation}

denote the ordinary unguided one-step ArchesWeather forecast for a fixed
generative noise realization \(\varepsilon_n\). The adversarial perturbation
is obtained from

\begin{equation}
    \delta_{n-1}^\star
    =
    \arg\min_{\delta\in\mathcal{C}}
    \left[
        \mathcal{M}_m
        \left(
            F_n(
                x_{n-2},
                x_{n-1}+\delta;
                \varepsilon_n
            )
        \right)
        -
        y_n^\star
    \right]^2 ,
    \label{eq:aowf-target-objective}
\end{equation}

where \(\mathcal{C}\) constrains the magnitude of the atmospheric perturbation.
The exact constraint, optimizer, and attack budget used in our experiments are
specified in the experimental setup.

Once \(\delta_{n-1}^\star\) has been obtained, generation itself remains
unguided. AOWF therefore differs fundamentally from the other methods in this
chapter: it modifies the atmospheric conditioning state before generation,
whereas the other approaches leave the atmospheric context fixed and intervene
inside the generative flow.

For comparison, AOWF is evaluated against the same prescribed weather-time
target trajectory \(y^\star_{1:N}\) and the same target intensities as the
guidance methods. Since the intervention occurs before the residual flow, no
flow-time guidance profile \(\gamma_t\) is associated with AOWF.

\section{Method comparison}
\label{sec:guidance-method-comparison}

The methods introduced above can be viewed as different answers to the same
target-realization problem. Their principal distinction is the location and
dimensionality of the intervention.

\begin{table}[t]
    \centering
    \caption{
        Conceptual comparison of the guidance and baseline methods.
        \(K\) denotes the number of outer optimization evaluations when
        applicable.
    }
    \label{tab:guidance-method-comparison}
    \small
    \begin{tabular}{p{0.18\linewidth} p{0.20\linewidth} p{0.22\linewidth} p{0.27\linewidth}}
        \toprule
        \textbf{Method}
        & \textbf{Intervention}
        & \textbf{Optimized quantity}
        & \textbf{Main characteristic} \\
        \midrule

        \textsc{Close-Gap}
        & Flow velocity
        & Local \(\lambda_{n,t}\)
        & Single-pass, locally follows a prescribed remaining-gap trajectory. \\

        \textsc{Optimize-Schedule}
        & Flow velocity
        & Full kick trajectory \(\mathbf{k}_n\)
        & Most flexible formulation; jointly optimizes interventions across flow time. \\

        \textsc{Optimize-Gain}
        & Flow velocity
        & Scalar \(w_n\)
        & Restricts guidance to \(w_n\gamma_t\), reducing target realization to a one-dimensional optimization. \\

        Universal Guidance
        & Clean estimate mapped to flow velocity
        & None
        & Constructs the target-closing clean-space correction directly and applies it according to \(\gamma_t\). \\

        AOWF
        & Atmospheric conditioning state
        & State perturbation \(\delta\)
        & Optimizes the model input and subsequently runs the generative forecast without flow guidance. \\

        \bottomrule
    \end{tabular}
\end{table}

The computational differences follow the same hierarchy.
\textsc{Close-Gap} and Universal Guidance determine their interventions online
and therefore require only a single guided residual flow per weather step.
\textsc{Optimize-Gain} requires repeated flow evaluations, but searches over
only one scalar parameter. \textsc{Optimize-Schedule} instead searches over
one intervention variable per guided flow step and requires repeated forward
and adjoint passes. AOWF also requires an outer optimization loop, but its
optimization is performed in atmospheric input space rather than over the
generative flow.

These differences motivate the experimental comparison in the following
chapters. Target realization alone is not sufficient to identify a useful
guidance method: methods must additionally be compared in terms of
meteorological plausibility, reliability across target intensities, sensitivity
to their hyperparameters, and computational cost.

Taken together, the methods in this chapter complete the \emph{Guide} component
of the proposed framework. Starting from the scenario specification
\((m,y^\star_{1:N})\), they provide different mechanisms for translating a
prescribed target into a modified generative forecast. The remaining question
is whether the resulting trajectories retain the structure of realistic
weather, which motivates the evaluation methodology introduced next.